\documentclass[11pt]{article}
\usepackage[utf8]{inputenc}
\usepackage{amsmath,amssymb,amsthm}
\usepackage[margin=1in]{geometry}
\usepackage{hyperref}
\usepackage{xcolor}
\usepackage{enumitem}
\usepackage{newunicodechar}
\newunicodechar{ℝ}{\ensuremath{\mathbb{R}}}
\newunicodechar{ι}{\ensuremath{\iota}}
\newunicodechar{φ}{\ensuremath{\varphi}}
\newunicodechar{ψ}{\ensuremath{\psi}}
\newunicodechar{→}{\ensuremath{\to}}
\newunicodechar{₁}{\ensuremath{_1}}
\newunicodechar{₂}{\ensuremath{_2}}
\newunicodechar{🟡}{\textbf{[in-progress]}}
\newunicodechar{✓}{\ensuremath{\checkmark}}

\title{Tide retrospective:\\multi-index affine-bilinear covariance mirror}
\author{Daniel Murfet}
\date{7 May 2026}

\begin{document}
\maketitle

\begin{abstract}
A 101-line mirror tide. The multi-index analogue of I3's affine-bilinear
covariance family,\footnote{The 1D namespace
\texttt{Laplace.gibbsCov\_affine\_*}; the multi-index mirror lives in
\texttt{Laplace.Multi}.} ported line-for-line to the new namespace.
Three public theorems (one-sided affine reductions on each slot, and
the headline two-sided identity); same four-witness integrability
budget; same proof skeleton; same \texttt{funext + ring} bridge for
the \texttt{Pi.add}-vs-single-lambda mismatch. The exercise is strictly
operational; the load-bearing observation is that the I3 abstraction
was at exactly the right level for the port to be mechanical. No
GPT-5.5~Pro consult; mirror-shape, proceeding-without-GPT path. 0
\texttt{sorry}.
\end{abstract}

\section{Setting}

I3 (the affine-bilinear specialisation tide, also 7~May) landed
three theorems in the 1D \texttt{gibbsCov} namespace demonstrating
the I2 algebra works in practice. The multi-index covariance on
$(\iota \to \mathbb{R})$ has the parallel-mirror algebra, also added
by I2; the natural next step is to mirror the I3 theorems into the
multi-index namespace, giving the affine-bilinear identity for both
1D and multi-index Gibbs measures.

The motivation is downstream: the 2D pure-quartic file currently
manually expands the affine-cov bashes that would collapse to
one-liners under the new helpers. The cleanup proper is the P1
follow-up; this tide just lands the algebra.

The candidates and proof skeleton were pre-determined by I3. Per the
tide skill's ``proceed-without-GPT'' path, established by the L1
sextic-J-function tide for mirror-shaped tides where the candidate is
already deliberated and the closed form is already verified, this
tide skips the GPT-5.5~Pro consult.

\section{The thing formalised}

\begin{quote}\itshape
For the multi-index Gibbs measure $\exp(-t \cdot L(w))\,\mathrm{d}w$
on $\iota \to \mathbb{R}$ (with $\iota$ a finite type) and any pair of
scalar observables $\varphi, \psi : (\iota \to \mathbb{R}) \to \mathbb{R}$
satisfying the four base integrability conditions, the multi-index
Gibbs covariance satisfies the affine-bilinear identity
\[
  \mathrm{Cov}_t[a_1 \varphi + b_1,\; a_2 \psi + b_2]
  \;=\; a_1 \cdot a_2 \cdot \mathrm{Cov}_t[\varphi, \psi]
  \qquad (a_1, b_1, a_2, b_2 \in \mathbb{R}).
\]
\end{quote}

The Lean signature of the headline\footnote{In
\texttt{Laplace/Multi/AffineBilinearCov.lean}.}, with $\iota$ a finite
type implicit:

\begin{verbatim}
theorem gibbsCov_affine_affine
    (L : (ι → ℝ) → ℝ) (t : ℝ) (a₁ b₁ a₂ b₂ : ℝ)
    (φ ψ : (ι → ℝ) → ℝ)
    (h_exp : Integrable (fun w => Real.exp (-(t * L w))))
    (h_φ   : Integrable (fun w => φ w * Real.exp (-(t * L w))))
    (h_ψ   : Integrable (fun w => ψ w * Real.exp (-(t * L w))))
    (h_φψ  : Integrable (fun w => φ w * ψ w * Real.exp (-(t * L w)))) :
    gibbsCov L t (fun w => a₁ * φ w + b₁) (fun w => a₂ * ψ w + b₂)
      = a₁ * a₂ * gibbsCov L t φ ψ
\end{verbatim}

The two one-sided helpers (\texttt{\_affine\_left} and
\texttt{\_affine\_right}) take the same four hypotheses and reduce a
single affine slot to a scalar multiple.

\section{Proof strategy}

Identical to I3's, with $w$ a function in $\iota \to \mathbb{R}$
in place of $x \in \mathbb{R}$ and the \texttt{Laplace.Multi.gibbsCov\_*}
algebra in place of the 1D version.

The single-slot reduction splits the affine combination via
\texttt{gibbsCov\_add\_left}, kills the constant via
\texttt{gibbsCov\_const\_left} and \texttt{add\_zero}, and factors the
scalar via \texttt{gibbsCov\_smul\_left}. Integrability witnesses
derive from the four base hypotheses by
\texttt{Integrable.const\_mul} and a \texttt{simpa [mul\_assoc]} cleanup.

The right-mirror is three lines via two applications of
\texttt{gibbsCov\_symm} sandwiching the left-version.

The headline composes the two one-sided versions with the right-slot
mixed integrability witnesses derived via \texttt{Integrable.add}
through a \texttt{funext + ring} bridge, and a final
\texttt{$\leftarrow$ mul\_assoc} to fold $a_1 \cdot (a_2 \cdot \cdot) = a_1 a_2 \cdot \cdot$.

\section{Roadblocks and resolutions}

None. The port was line-for-line identical to I3, except for
the lambda variable name ($w$ instead of $x$) and the namespace
prefix. The expected friction points (the \texttt{Pi.add}-vs-single-lambda
mismatch, the \texttt{simpa [mul\_assoc]} cleanup, the
\texttt{Integrable.const\_mul + Integrable.add} construction) all
worked verbatim from the I3 template.

The exercise lands first try at every step. \texttt{lake build}
succeeded on the first compile.

The 101-line file size matches I3's 101-line size to the line, an
unexpectedly precise match; structurally the two files are
isomorphic.

\section{What was Mathlib, what was new}

\paragraph{Mathlib.} \texttt{Integrable.const\_mul},
\texttt{Integrable.add}. Standard ring tactics.

\paragraph{New.} Three theorems and zero infrastructure
(\texttt{gibbsCov\_affine\_left/right/affine\_affine} in
\texttt{Laplace.Multi}). Everything else was already present from I2's
multi-index algebra, the multi-index Gibbs definitions, and I3's
proof structure.

\section{Lessons}

\begin{enumerate}[leftmargin=*]

\item \textbf{Mirror tides are nearly free when I2-shaped abstractions
exist.} The I2 algebra was built explicitly to factor across the 1D
and multi-index Gibbs definitions, with parallel-mirror lemma names
(\texttt{Laplace.gibbsCov\_*} and \texttt{Laplace.Multi.gibbsCov\_*}).
Once I3 worked in 1D, the multi-index mirror was three minutes after
the tide-log was drafted, plus build verification.

\item \textbf{Mirror-shape tides should skip the GPT consult.} The L1
sextic-J-function tide established the precedent; this tide
re-confirms it. When the candidate is identical-modulo-namespace to a
just-landed tide, the deliberation has already happened in the
predecessor; the \texttt{tide-pick} \texttt{🟡 In progress} → \texttt{✓
Closed} cycle still applies, but the Step 2 GPT round is overhead.

\item \textbf{The right time to land a mirror tide is right after the
predecessor.} I3 landed earlier this session with the proof template
fresh in cache; mirroring it now meant zero context-loading cost. The
same template a week from now would require re-reading I3's
retrospective and the laplace \texttt{CLAUDE.md} to recover the
\texttt{Pi.add}-vs-single-lambda fix. Strike while the iron is hot.

\end{enumerate}

\section{Follow-ups}

\begin{itemize}[leftmargin=*]

\item \textbf{2D pure-quartic affine-cov cleanup} (the natural P1
sub-tide that this tide unblocks). Now that
\texttt{Laplace.Multi.gibbsCov\_affine\_*} exists, the
2D pure-quartic file's affine-cov bashes can be replaced with
one-liners. Estimated $\sim$80 line net deletion. Should be a separate
tide (separate retrospective; cleaner narrative).

\item \textbf{Promote the affine helpers to the
\texttt{gibbsCov}-algebra blocks.} Both the 1D and multi-index
versions now exist as sibling files; after 2--3 consumer sites land
for each, fold the helpers into the
\texttt{gibbsCov}-algebra sections of the corresponding base Gibbs
files so they sit alongside their predecessors.

\item \textbf{Threepoint-side analogue (Q4 in the v4 survey).}
Threepoint's \texttt{gibbsCov} (with $h \cdot A$ in the potential) has
no algebra at all. A port of I2 + I3 to threepoint would unblock
$\kappa_3$ strict-improvements stacked on the perturbed measure.
Bigger; cross-repo.

\end{itemize}

\end{document}
